\documentclass{article}
\usepackage{amsmath}
\usepackage{xcolor}
\begin{document}

\textbf{APLICACIONES DE MODELADO PARA E.D. de Primer Orden}



Aplicaciones de Crecimiento/Decrecimiento(Decaimiento)

Dada una determinado variable, se requiere determinar su comportamiento a lo largo del tiempo

Cuando la cantidad presente de la variable de estudio es directamente proporcional a una ctte,  se

puede representar la misma de la siguiente forma, a traves de una E.D.


\begin{gather*}
\frac{dP}{dt} = kP ;  P(t_{0}) = P_{0}\text{ donde k} \rightarrow \text{ ctte de proporcionalidad}
\end{gather*}


La E.D es de variable separable y posible resolverla


\begin{gather*}
\frac{1}{P}dP = kdt  \rightarrow  \int \frac{1}{P}dP = k\int dt \\
ln(P) = kt+c \\
P(t) = e^{kt+c}=Ae^{kt}, donde A=e^{c}
\end{gather*}


\textbf{Ejercicios}

Inicialmente un cultivo tiene un número 
\begin{gather*}
P_{0}
\end{gather*}
 de bacterias. En t = 1 h se determina que

el número de bacterias es 
\begin{gather*}
\frac{3}{2}P_{0}
\end{gather*}
. Si la razón de crecimiento es proporcional al número

de bacterias P(t) presentes en el tiempo t, determine el tiempo necesario para que se

triplique el número de bacterias.

\image-container


\begin{gather*}
P(t) =Ae^{kt}
\end{gather*}


Se toma inicialmente un punto de estudio

Tomando como punto de partida el \colorbox[RGB]{248,231,28}{punto A}


\begin{gather*}
P_{0} =Ae^{k(0)} \rightarrow    A=P_{0} \\
 P(t) =P_{0}e^{kt} 
\end{gather*}


Ahora tomamos el siguiente\colorbox[RGB]{248,231,28}{ punto B}


\begin{gather*}
\frac{3P_{0}}{2} =P_{0}e^{k(1)}  \rightarrow \frac{3\text{\colorbox[RGB]{248,231,28}{$P$}}\text{\colorbox[RGB]{248,231,28}{$_{0}$}}}{2\text{\colorbox[RGB]{248,231,28}{$P$}}\text{\colorbox[RGB]{248,231,28}{$_{0}$}}}=e^{k}  \rightarrow k=ln(3/2) \\
k=0.41 \\
P(t) =P_{0}e^{0.41t} 
\end{gather*}


Para el punto C, se nos pide el tiempo


\begin{gather*}
3P_{0} =P_{0}e^{0.41t} \rightarrow  t=\frac{ln(3)}{0.41}  
\end{gather*}
t = 2.679



Si hubiesen pasado 5hrs, cuanto de poblacion habria?


\begin{gather*}
P(_{t=5})= 
\end{gather*}
 
\begin{gather*}
P_{0}e^{0.41(5)} =P_{0}(2.71..) 
\end{gather*}




\textbf{LEY DE NEWTON, enfriamiento y Calentamiento de Cuerpos}


\begin{gather*}
\frac{dT}{dt}= k(T - T_{m})
\end{gather*}



\begin{gather*}
\int \frac{dT}{(T - T_{m})}= k\int dt \\
ln(T - T_{m}) =kt+c \\
T-T_{m}=e^{kt+c}=Ae^{kt} \\
T=Ae^{kt}+T_{m}
\end{gather*}
    



Una taza de cafe fresco se ha preparado en la cafetera y actualmente tiene una Temperatura de 200ºF. Despues de 3 minutos, este se ha enfriado y ha alcanzado una temperatura de 190ºF, pero aun esta un tanto caliente para ser consumido. Si la temperatura del ambiente es de 72ºF y la temperatura del Cafe esta regido a la ley de Enfriamiento de Newton, Que tiempo sera necesario para que se pueda consumir la taza de cafe de tal manera que haya alcanzado una temperatura de 165F?\image-container


\begin{gather*}
\text{Inicialmente aplicamos la ley de Newton:}
\end{gather*}



\begin{gather*}
T=Ae^{kt}+T_{m} \\
\text{Primeramente arrancamos con }(A) \\
T(0) =200  \rightarrow \text{ la termperatura del cafe arranca con ese valor} \\
200 = Ae^{k(0)}+T_{m} \\
200 = A+72 \\
A = 200-72 \\
A=128 \\
\text{Reemplazamos en la expresion inicial} \\
T=128e^{kt}+72 \\
\vspace{0.5em} \\
\text{Continuamos con el  punto }(B) \\
T(3) = 190 \\
T=128e^{kt}+72 \\
190 = 128e^{k(3)}+72 \\
\frac{ln(\frac{190-72}{128})}{3}= k  \rightarrow  k =-0.027 \\
\text{Reemplazamos en la expresion inicial} \\
T=128e^{-0.027t}+72 \\
\vspace{0.5em} \\
\text{Para determinar el valor del tiempo, reemplazamos en la formula anterior} \\
T(t=?) = 165 \\
165=128e^{-0.027t}+72 \\
\frac{165-72}{128} =e^{-0.027t} \\
t=\frac{ln(\frac{165-72}{128})}{-0.027} =11.83s
\end{gather*}




\image-container





Supon que un virus se esta esparciendo por toda la poblacion a un valor directamente proporcional al numero de personas que no estan enfermas. Si el 5% de la poblacion inicialmente se encuentra ya infectada y posterior a una semana se enferma ya el 10% de la poblacion. Encuentra la formula como tal que representa el modelo de la poblacion infectada en cualquier tiempo.


\begin{gather*}
\text{La representacion del modelo estara dado por : Asumimos que el universo de personas } \\
\text{sanas estara representada por el 100% } \rightarrow  1
\end{gather*}





\begin{gather*}
\frac{dP}{dt} = k(1-P) \\
\text{Resolvamos la E.D.} \\
\int \frac{1}{1-P}dp = \int kdt \\
-ln(1-P) = kt+c \\
ln(1-P)^{-1} = kt+c \\
\frac{1}{1-P}=Ae^{kt} \\
1-P=Ae^{-kt} \\
P=1-Ae^{-kt}
\end{gather*}




\image-container




\begin{gather*}
\text{Para ambos Puntos} \\
 P(0) = 0.05   \rightarrow (A) \\
P(1) = 0.1  \rightarrow  ( B) \\
\text{Con los valores de frontera es posible determinar el modelo de la E.D.} \\
0.05=1-Ae^{-k(0)} \\
0.05=1-A \\
A=0.95 \\
\vspace{0.5em} \\
\text{Para el siguiente punto Frontera:} \\
P=1-0.95e^{-kt} \\
0.1=1-0.95e^{-k(1)} \\
0.1-1=-0.95e^{-k(1)} \\
-0.9=-0.95e^{-k(1)} \\
\frac{0.9}{0.95} =e^{-k}  \rightarrow k=-ln(\frac{0.9}{0.95} )=0.054 \\
P=1-0.95e^{-0.054t}
\end{gather*}





\end{document}
